\section{AMC 12 12A 2010}
        \begin{enumerate}
        	\item (\textit{AMC 12 12A 2010}) What is $\left(20-\left(2010-201\right)\right)+\left(2010-\left(201-20\right)\right)$?


 $\textbf{(A)}\ -4020 \qquad \textbf{(B)}\ 0 \qquad \textbf{(C)}\ 40 \qquad \textbf{(D)}\ 401 \qquad \textbf{(E)}\ 4020$

 \textbf{Answer: }C



\textbf{Solution}

$20-2010+201+2010-201+20=20+20=\boxed{\textbf{(C)}\,40}$.

 


	\item (\textit{AMC 12 12A 2010}) A ferry boat shuttles tourists to an island every hour starting at 10 AM until its last trip, which starts at 3 PM. One day the boat captain notes that on the 10 AM trip there were 100 tourists on the ferry boat, and that on each successive trip, the number of tourists was 1 fewer than on the previous trip. How many tourists did the ferry take to the island that day?


 $\textbf{(A)}\ 585 \qquad \textbf{(B)}\ 594 \qquad \textbf{(C)}\ 672 \qquad \textbf{(D)}\ 679 \qquad \textbf{(E)}\ 694$

 \textbf{Answer: }A



\textbf{Solution 1}

It is easy to see that the ferry boat takes $6$ trips total. The total number of people taken to the island is

 
 \begin{align*}&100+99+98+97+96+95\\  &=6(100)-(1+2+3+4+5)\\  &=600 - 15\\  &=\boxed{585\ \textbf{(A)}}\end{align*}





\textbf{Solution 2}

You can notice that:
 \begin{align*}&100 + 95 = 99 + 96 = 98 + 97 = 195\end{align*}
This means that the sum is:
 \begin{align*}&100+99+98+97+96+95\\  &=195 * 3\\  &=\boxed{585\ \textbf{(A)}}\end{align*}




	\item (\textit{AMC 12 12A 2010}) Rectangle $ABCD$, pictured below, shares $50\%$ of its area with square $EFGH$. Square $EFGH$ shares $20\%$ of its area with rectangle $ABCD$. What is $\frac{AB}{AD}$?


 
                    \begin{figure}[H]
                        \centering
                        \includegraphics[width=0.5\textwidth]{Problem/AMC_12_12A_2010_3_Problem_Figure_1.png}
                    \end{figure}
                    
$\textbf{(A)}\ 4 \qquad \textbf{(B)}\ 5 \qquad \textbf{(C)}\ 6 \qquad \textbf{(D)}\ 8 \qquad \textbf{(E)}\ 10$

 \textbf{Answer: }E



\textbf{Solution 1}

If we shift $A$ to coincide with $E$, and add new horizontal lines to divide $EFGH$ into five equal parts:

 
                    \begin{figure}[H]
                        \centering
                        \includegraphics[width=0.5\textwidth]{Solution/AMC_12_12A_2010_3_Solution_Figure_1.png}
                    \end{figure}
                    

 This helps us to see that $AD=a/5$ and $AB=2a$, where $a=EF$.Hence $\dfrac{AB}{AD}=\dfrac{2a}{a/5}=10$.

 



\textbf{Solution 2}

From the problem statement, we know that
 \[\frac{[ABCD]}{2} = \frac{[EFGH]}{5} \Rightarrow [ABCD]=\frac{2[EFGH]}{5}\]

 If we let $a = EF$ and $b  = AD$, we see
 \[[ABCD] = 2ab = \frac{2a^2}{5} \Rightarrow b = \frac{a}{5}\]. Hence, $\frac{AB}{AD} = \frac{2a}{b} = 2a(\frac{5}{a}) = \boxed{\bold{10}}$



	\item (\textit{AMC 12 12A 2010}) If $x<0$, then which of the following must be positive?


 $\textbf{(A)}\ \frac{x}{\left|x\right|} \qquad \textbf{(B)}\ -x^2 \qquad \textbf{(C)}\ -2^x \qquad \textbf{(D)}\ -x^{-1} \qquad \textbf{(E)}\ \sqrt[3]{x}$

 \textbf{Answer: }D



\textbf{Solution}

$x$ is negative, so we can just place a negative value into each expression and find the one that is positive. Suppose we use $-1$.

 $\textbf{(A)} \Rightarrow \frac{-1}{|-1|} = -1$
$\textbf{(B)} \Rightarrow -(-1)^2 = -1$
$\textbf{(C)} \Rightarrow -2^{(-1)} = -\frac{1}{2}$
$\textbf{(D)} \Rightarrow -(-1)^{(-1)} = 1$
$\textbf{(E)} \Rightarrow \sqrt[3]{-1} = -1$

 
Obviously only $\boxed{\textbf{(D)}}$ is positive.

 


	\item (\textit{AMC 12 12A 2010}) Halfway through a 100-shot archery tournament, Chelsea leads by 50 points. For each shot a bullseye scores 10 points, with other possible scores being 8, 4, 2, and 0 points. Chelsea always scores at least 4 points on each shot. If Chelsea's next $n$ shots are bullseyes she will be guaranteed victory. What is the minimum value for $n$?


 $\textbf{(A)}\ 38 \qquad \textbf{(B)}\ 40 \qquad \textbf{(C)}\ 42 \qquad \textbf{(D)}\ 44 \qquad \textbf{(E)}\ 46$

 \textbf{Answer: }C



\textbf{Solution}

Let $k$ be the number of points Chelsea currently has. In order to guarantee victory, we must consider the possibility that the opponent scores the maximum amount of points by getting only bullseyes.

 
 \begin{align*}k+ 10n + 4(50-n) &> (k-50) + 50\cdot{10}\\  6n &> 250\end{align*}


 
The lowest integer value that satisfies the inequality is $\boxed{42\ \textbf{(C)}}$.

 


	\item (\textit{AMC 12 12A 2010}) A $\textit{palindrome}$, such as 83438, is a number that remains the same when its digits are reversed. The numbers $x$ and $x+32$ are three-digit and four-digit palindromes, respectively. What is the sum of the digits of $x$?


 $\textbf{(A)}\ 20 \qquad \textbf{(B)}\ 21 \qquad \textbf{(C)}\ 22 \qquad \textbf{(D)}\ 23 \qquad \textbf{(E)}\ 24$

 \textbf{Answer: }E



\textbf{Solution 1}

$x$ is at most $999$, so $x+32$ is at most $1031$. The minimum value of $x+32$ is $1000$. However, the only palindrome between $1000$ and $1032$ is $1001$, which means that $x+32$ must be $1001$.

 It follows that $x$ is $969$, so the sum of the digits is $\boxed{\textbf{(E)}\ 24}$.

 



\textbf{Solution 2}

For $x+32$ to be a four-digit number, $x$ is in between $968$ and $999$. The palindromes in this range are $969$, $979$, $989$, and $999$, so the sum of the digits of $x$ can be $24$, $25$, $26$, or $27$. Only $\boxed{\textbf{(E)}\ 24}$ is an option, and upon checking, $x+32=1001$ is indeed a palindrome.

 



\textbf{Solution 3}

Since we know $x+32$ to be $1 a a 1$ and the only palindrome that works is $0 = a$, that means $x+32 = 1001$, and so $x = 1001 - 32 = 969$. So, $9$ + $6$ + $9$ = $\boxed{\textbf{(E)}\ 24}$.~songmath20

 


	\item (\textit{AMC 12 12A 2010}) Logan is constructing a scaled model of his town. The city's water tower stands 40 meters high, and the top portion is a sphere that holds 100,000 liters of water. Logan's miniature water tower holds 0.1 liters. How tall, in meters, should Logan make his tower?


 $\textbf{(A)}\ 0.04 \qquad \textbf{(B)}\ \frac{0.4}{\pi} \qquad \textbf{(C)}\ 0.4 \qquad \textbf{(D)}\ \frac{4}{\pi} \qquad \textbf{(E)}\ 4$

 \textbf{Answer: }C



\textbf{Solution}

The water tower holds $\frac{100,000}{0.1} = 1,000,000$ times more water than Logan's miniature. The volume of a sphere is: $V=\dfrac{4}{3}\pi r^3$. Since we are comparing the heights (m), we should compare the radii (m) to find the ratio. Since, the radius is cubed, Logan should make his tower $\sqrt[3]{1,000,000} = 100$ times shorter than the actual tower. This is $\frac{40}{100} = \boxed{0.4}$ meters high, or choice $\textbf{(C)}$.

 Note: The fact that $1\text{ L}=1,000\text{ cm}^3$ doesn't matter since only the ratios are important.

 


	\item (\textit{AMC 12 12A 2010}) Triangle $ABC$ has $AB=2 \cdot AC$. Let $D$ and $E$ be on $\overline{AB}$ and $\overline{BC}$, respectively, such that $\angle BAE = \angle ACD$. Let $F$ be the intersection of segments $AE$ and $CD$, and suppose that $\triangle CFE$ is equilateral. What is $\angle ACB$?


 $\textbf{(A)}\ 60^\circ \qquad \textbf{(B)}\ 75^\circ \qquad \textbf{(C)}\ 90^\circ \qquad \textbf{(D)}\ 105^\circ \qquad \textbf{(E)}\ 120^\circ$

 \textbf{Answer: }C



\textbf{Solution 1}


                    \begin{figure}[H]
                        \centering
                        \includegraphics[width=0.5\textwidth]{Solution/AMC_12_12A_2010_8_Solution_Figure_1.png}
                    \end{figure}
                    

 
Let $\angle BAE = \angle ACD = x$.

 
 \begin{align*}\angle BCD &= \angle AEC = 60^\circ\\  \angle EAC + \angle FCA + \angle ECF + \angle AEC &= \angle EAC + x + 60^\circ + 60^\circ = 180^\circ\\  \angle EAC &= 60^\circ - x\\  \angle BAC &= \angle EAC + \angle BAE = 60^\circ - x + x = 60^\circ\end{align*}


 Since $\frac{AC}{AB} = \frac{1}{2}$ and the angle between the hypotenuse and the shorter side is $60^\circ$, triangle $ABC$ is a $30-60-90$ triangle, so $\angle BCA = \boxed{90^\circ\,\textbf{(C)}}$.

 



\textbf{Solution 2(Trig and Angle Chasing)}

Let 
 \[AB=2a, AC=a\]Let 
 \[\angle BAE=\angle ACD=x\]Because $\triangle CFE$ is equilateral, we get $\angle FCE=60$, so $\angle ACB=60+x$

 Because $\triangle CFE$ is equilateral, we get $\angle CFE=60$. 

 Angles $AFD$ and $CFE$ are vertical, so $\angle AFD=60$. 

 By triangle $ADF$, we have $\angle ADF=120-x$, and because of line $AB$, we have $\angle BDC=60+x$.

 Because Of line $BC$, we have $\angle AEB=120$, and by line $CD$, we have $\angle DFE=120$. 

 By quadrilateral $BDFE$, we have $\angle ABC=60-x$.

 By the Law of Sines:
 \[\frac{\sin(60-x)}{a}=\frac{\sin(60+x)}{2a}\implies 2\sin(60-x)=\sin(60+x)\]

 By the sine addition formula($\sin(a+b)=\sin(a)\cos(b)+\cos(a)\sin(b)$):
 \[2(\sin(60)\cos(-x)+\cos(60)\sin(-x))=\sin(60)\cos(x)+\cos(60)\sin(x)\]

 Because cosine is an even function, and sine is an odd function, we have 
 \[2\sin(60)\cos(x)-2\cos(60)\sin(x)=\sin(60)\cos(x)+\cos(60)\sin(x) \implies \sin(60)\cos(x)=3\cos(60)\sin(x)\]

 We know that $\sin(60)=\frac{\sqrt{3}}{2}$, and $\cos(60)=\frac{1}{2}$, hence 
 \[\frac{\sqrt{3}}{2}\cos(x)=\frac{3}{2}\sin(x)\implies \tan(x)=\frac{\sqrt{3}}{3}\]

 The only value of $x$ that satisfies $60+x<180$(because $60+x$ is an angle of the triangle) is $x=30^{\circ}$. We seek to find $\angle ACB$, which as we found before is $60+x$, which is $90$. The answer is $90, \text{or} \textbf{(C)}$

 -vsamc

 



\textbf{Solution 3 (Similar Triangles)}

Notice that $\angle AEB=\angle AFC = 120^{\circ}$ and $\angle ACF=\angle BAE$. Hence, triangle AEB is similar to triangle CFA. Since $AB=2AC$, $AE=2CF=2FE$, as triangle CFE is equilateral. Therefore, $AF=FE=FC$, and since $\angle AFC=120^{\circ}$, $x=30$. Thus, the measure of $\angle ACE$ equals to $\angle FCE+\angle ACF=90^{\circ}, \text{or} \textbf{(C)}$ -HarryW

 


 



\textbf{Solution 4}

Notice that $\triangle ADF \sim \triangle CDA$ (by AA Similarity.) Since the corresponding angles of a pair of similar triangles are congruent, we have $\angle DAC = \angle DFA = 60^\circ.$ Since $AB = 2 \cdot AC$ and $\angle DFA = 60^\circ,$ we have that $\triangle ABC$ is congruent by SAS to a $30-60-90$ right triangle, which gives the answer $\boxed{\textbf{(C)}}$.~clever14710owl

 


	\item (\textit{AMC 12 12A 2010}) A solid cube has side length $3$ inches. A $2$-inch by $2$-inch square hole is cut into the center of each face. The edges of each cut are parallel to the edges of the cube, and each hole goes all the way through the cube. What is the volume, in cubic inches, of the remaining solid?


 $\textbf{(A)}\ 7 \qquad \textbf{(B)}\ 8 \qquad \textbf{(C)}\ 10 \qquad \textbf{(D)}\ 12 \qquad \textbf{(E)}\ 15$

 \textbf{Answer: }A



\textbf{Solution 1}

Imagine making the cuts one at a time. The first cut removes a box $2\times 2\times 3$. The second cut removes two boxes, each of dimensions $2\times 2\times 0.5$, and the third cut does the same as the second cut, on the last two faces. Hence the total volume of all cuts is $12 + 4 + 4 = 20$.

 Therefore the volume of the rest of the cube is $3^3 - 20 = 27 - 20 = \boxed{7\ \textbf{(A)}}$.

 



\textbf{Solution 2}

We can visualize it has a $3$ by $3$ cube with a hollow $2$ by $2$ center, along with 6 "windows" cut out too. The cube formula gives the volume of the $3$ by $3$ cube as $3^3 = 27$ and $2^3 = 8.$ Additionally, the area of the "windows" is the 2x2 cut along with a thickness of $\frac{(3-2)}{2} = \frac{1}{2}.$ The volume of all 6 windows is $6(\frac{1}{2})(2)(2) = 12.$ This gives us our answer of $27 - 8 - 12 = \boxed{7\ \textbf{(A)}}.$




\textbf{Solution 3}

We can use \href{https://artofproblemsolving.com/wiki/index.php?title=Principle\_of\_Inclusion-Exclusion}{Principle of Inclusion-Exclusion} (PIE) to find the final volume of the cube.

 There are 3 "cuts" through the cube that go from one end to the other. Each of these "cuts" has $2 \times 2 \times 3=12$ cubic inches. However, we can not just sum their volumes, asthe central $2\times 2\times 2$ cube is included in each of these three cuts. To get the correct result, we can take the sum of the volumes of the three cuts, and subtract the volume of the central cube twice.

 Hence the total volume of the cuts is $3(2 \times 2 \times 3) - 2(2\times 2\times 2) = 36 - 16 = 20$.

 Therefore the volume of the rest of the cube is $3^3 - 20 = 27 - 20 = \boxed{7\ \textbf{(A)}}$.

 



\textbf{Solution 4}

We can visualize the final figure and see a cubic frame. We can find the volume of the figure by adding up the volumes of the edges and corners.

 Each edge can be seen as a $2\times 0.5\times 0.5$ box, and each corner can be seen as a $0.5\times 0.5\times 0.5$ box.

 $12\cdot{\frac{1}{2}} + 8\cdot{\frac{1}{8}} = 6+1 = \boxed{7\ \textbf{(A)}}$.

 


	\item (\textit{AMC 12 12A 2010}) The first four terms of an arithmetic sequence are $p$, $9$, $3p-q$, and $3p+q$. What is the $2010^\text{th}$ term of this sequence?


 $\textbf{(A)}\ 8041 \qquad \textbf{(B)}\ 8043 \qquad \textbf{(C)}\ 8045 \qquad \textbf{(D)}\ 8047 \qquad \textbf{(E)}\ 8049$

 \textbf{Answer: }A



\textbf{Solution 1}

$3p-q$ and $3p+q$ are consecutive terms, so the common difference is $(3p+q)-(3p-q) = 2q$.

 
 \begin{align*}p+2q &= 9\\ 9+2q &= 3p-q\\ q&=2\\ p&=5\end{align*}


 The common difference is $4$. The first term is $5$ and the $2010^\text{th}$ term is

 
 \[5+4(2009) = \boxed{\textbf{(A) }8041}\]




\textbf{Solution 2}

Since all the answer choices are around $2010 \cdot 4 = 8040$, the common difference must be $4$. The first term is therefore $9 - 4 = 5$, so the $2010^\text{th}$ term is $5 + 4 \cdot 2009 = \boxed{\textbf{(A) }8041}$.

 


	\item (\textit{AMC 12 12A 2010}) The solution of the equation $7^{x+7} = 8^x$ can be expressed in the form $x = \log_b 7^7$. What is $b$?


 $\textbf{(A)}\ \frac{7}{15} \qquad \textbf{(B)}\ \frac{7}{8} \qquad \textbf{(C)}\ \frac{8}{7} \qquad \textbf{(D)}\ \frac{15}{8} \qquad \textbf{(E)}\ \frac{15}{7}$

 \textbf{Answer: }C



\textbf{Solution 1}

This problem is quickly solved with knowledge of the laws of exponents and logarithms.

 
 \begin{align*} 7^{x+7} &= 8^x \\  7^x*7^7 &= 8^x \\  \left(\frac{8}{7}\right)^x &= 7^7 \\  x &= \log_{8/7}7^7 \end{align*}


 Since we are looking for the base of the logarithm, our answer is $\boxed{\textbf{(C)}\ \frac{8}{7}}$.

 



\textbf{Solution 2}

First, take the $\log_{7}$ of both sides, which gives us $x+7=\log_{7}8^x=x\log_{7}8$. We move the $x$ terms to 1 side. $x(\log_{7}8-1)=7$. Isolate $x$ and manipulate the answer. $x=\frac{7}{\log_{7}\frac{8}{7}}=7\log_{\frac{8}{7}}7=\log_{\frac{8}{7}}7^7$. Therefore, the answer is $\frac{8}{7}=\fbox{C}$



	\item (\textit{AMC 12 12A 2010}) In a magical swamp there are two species of talking amphibians: toads, whose statements are always true, and frogs, whose statements are always false. Four amphibians, Brian, Chris, LeRoy, and Mike live together in this swamp, and they make the following statements.


 Brian: "Mike and I are different species."


 Chris: "LeRoy is a frog."


 LeRoy: "Chris is a frog."


 Mike: "Of the four of us, at least two are toads."


 How many of these amphibians are frogs?


 $\textbf{(A)}\ 0 \qquad \textbf{(B)}\ 1 \qquad \textbf{(C)}\ 2 \qquad \textbf{(D)}\ 3 \qquad \textbf{(E)}\ 4$

 \textbf{Answer: }D



\textbf{Solution 1}

Start with Brian. If he is a toad, he tells the truth, hence Mike is a frog. If Brian is a frog, he lies, hence Mike is a frog, too. Thus Mike must be a frog.

 As Mike is a frog, his statement is false, hence there is at most one toad.

 As there is at most one toad, at least one of Chris and LeRoy is a frog. But then the other one tells the truth, and therefore is a toad.

 Hence we must have one toad and $\boxed{\textbf{(D)}\ 3}$ frogs.

 



\textbf{Solution 2 (logical reasoning like solution 1, but a different train of thought)}

Notice that one of Chris and LeRoy must be a frog: if Chris is a frog, then he lies about LeRoy being a frog. Hence LeRoy is a toad. Alternatively, if Chris is a toad, then he tells the truth about LeRoy being a frog.

 Assume Brian is a toad. Then Mike is a frog, and he lies about at least two being toads. This means that none or one of the amphibians is a toad (the opposite of the statement $n\geq2$ is $n<2$, or $n=0, 1$). However, this is absurd because we assumed Brian is a toad, and we know one of Chris and LeRoy is a toad. So our assumption leads to a contradiction.

 Hence Brian must be a frog, and he and Mike are the same species. Mike is also a frog. One of Chris and LeRoy is a frog. There are $3$ frogs in total $\Longrightarrow \boxed{\textbf{(D) } 3}$.

 



\textbf{Solution 3: Casework/Proof by contradiction}

First, we assume that Brian is a toad. Since his statements are truthful, Mike could not be a toad and must be a frog. Since frogs are liars, Brian must be the only toad in the pond, so Chris and Leroy are both frogs. However, they would then be telling the truth about each other, which only toads do. Hence, this case results in a contraction.Hence, Brian is a frog. He is lying, so Mike must also be a frog. Mike's statement is also false, so Chris and Leroy cannot both be toads. As shown in the previous case, them both being frogs is impossible, so one must be a toad and the other a frog. Therefore, there must be a total of 3 frogs.By-Monkey\_King

 


	\item (\textit{AMC 12 12A 2010}) For how many integer values of $k$ do the graphs of $x^2+y^2=k^2$ and $xy = k$ not intersect?


 $\textbf{(A)}\ 0 \qquad \textbf{(B)}\ 1 \qquad \textbf{(C)}\ 2 \qquad \textbf{(D)}\ 4 \qquad \textbf{(E)}\ 8$

 \textbf{Answer: }C



\textbf{Solution 1}

The image below shows the two curves for $k=4$. The blue curve is $x^2+y^2=k^2$, which is clearly a circle with radius $k$, and the red curve is a part of the curve $xy=k$.

 
                    \begin{figure}[H]
                        \centering
                        \includegraphics[width=0.5\textwidth]{Solution/AMC_12_12A_2010_13_Solution_Figure_1.png}
                    \end{figure}
                    

 In the special case $k=0$ the blue curve is just the point $(0,0)$, and as $0\cdot 0=0$, this point is on the red curve as well, hence they intersect. 

 The case $k<0$ is symmetric to $k>0$: the blue curve remains the same and the red curve is flipped according to the $x$ axis. Hence we just need to focus on $k>0$.

 Clearly, on the red curve there will always be points arbitrarily far from the origin: for example, as $x$ approaches 0, $y$ approaches $\infty$. Hence the red curve intersects the blue one if and only if it contains a point whose distance from the origin is at most $k$.

 
At this point we can guess that on the red curve the point where $x=y$ is always closest to the origin, and skip the rest of this solution.

 
For an exact solution, fix $k$ and consider any point $(x,y)$ on the red curve. Its distance from the origin is $\sqrt{ x^2 + (k/x)^2 }$. To minimize this distance, it is enough to minimize $x^2 + (k/x)^2$. By the \href{https://artofproblemsolving.com/wiki/index.php?title=Arithmetic\_Mean-Geometric\_Mean\_Inequality}{Arithmetic Mean-Geometric Mean Inequality} we get that this value is at least $2k$, and that equality holds whenever $x^2 = (k/x)^2$, i.e., $x=\pm\sqrt k$.

 
Now recall that the red curve intersects the blue one if and only if its closest point is at most $k$ from the origin. We just computed that the distance between the origin and the closest point on the red curve is $\sqrt{2k}$. Therefore, we want to find all positive integers $k$ such that $\sqrt{2k} > k$.

 Clearly the only such integer is $k=1$, hence the two curves are only disjoint for $k=1$ and $k=-1$. This is a total of $\boxed{2\ \textbf{(C)}}$ values.

 



\textbf{Solution 2}

From the graph shown above, we see that there is a specific point closest to the center of the circle. Using some logic, we realize that as long as said furthest point is not inside or on the graph of the circle. This should be enough to conclude that the hyperbola does not intersect the circle. 

 Therefore, for each value of k, we only need to check said value to determine intersection. Let said point, closest to the circle have coordinates $(x, k/x)$ derived from the equation. Then, all coordinates that satisfy $\sqrt{ x^2+ (k/x)^2 } \leq k$ intersect the circle.Squaring, we find $x^2+(k/x)^2 \leq k^2.$After multiplying through by $x^2$ and rearranging, we find $x^4-x^2k^2+k^2 \leq 0$.We see this is a quadratic in $x^2$ and consider taking the determinant, which tells us that solutions are real when, after factoring:$k^2(k^2-4) \geq 0$We plot this inequality on the number line to find it is satisfied for all values except: $(-1, 0, 1)$.We then eliminate 0 because it is extraneous as both $xy=0$ and $x^2+y^2=0$ are points which coincide.Therefore, there are a total of $\boxed{2\ \textbf{(C)}}$ values.

 



\textbf{Solution 3 (Algebra)}

Since $xy=k$, multiply the equation by 2 on both sides to get $2xy=2k$. Now we can add the two equations to get $(x+y)^2=k^2+2k$, for which the only value of $k$ that does not satisfy the equation is $-1$, as that makes the RHS negative. Similarly, if we subtract the two equations, we obtain $(x-y)^2=k^2-2k$, for which the only value of $k$ that does not satisfy the equation is $1$, for the same reason above.

 Thus, the only values are $k = 1, -1$, giving us a total of $\boxed{2\ \textbf{(C)}}$ values.

 



\textbf{Solution 4 (Quick)}

Multiply $k=xy$ by $k$ and substitute it into $k^2=x^2+y^2$. Then, $k=\frac{x^2+y^2}{xy}$. Recognize it? It's also $k=\frac{x}{y}+\frac{y}{x}$. The minimum of this function (more accurately the minimum absolute value of the function) is $k=2, -2$ (when $x=y$ or $x=-y$). As long as $k>2$ or $k<-2$, the function is valid. As such, $k\neq1,-1 \implies \boxed{2\ \textbf{(C)}}$. Elegant, huh?

 



\textbf{Solution 5}

Assume that $k\ge 0$ since if $k$ works then $-k$ also works. Let $x = ka$ and $y = kb$. Then the given equations become $ab = \frac{1}{k}$ and $a^2 + b^2 = 1$, which we don't want to intersect. The points that are the closest on this graph are $(1/\sqrt{2}, 1/\sqrt{2})$ and $(1/\sqrt{k})(1/\sqrt{k})$. We don't want the hyperbola to intersect or go inside the circle, so we require $1/\sqrt{k} > 1/\sqrt{2}$, so $k< 2\implies k = 0, 1$. Also, by symmetry $k = -1$ also works.Obviously $k = 0$ doesn't work, so we can discard that, leaving $\boxed{\mathbf{(C)}\quad 2}$ integral solutions for $k$




\textbf{Solution 6(Algebra)}

To know when the system does not intersect, we first need to find when they do intersect. To do so, we substitute $xy=k$ into $x^{2}+y^{2}=k^{2}$ giving:

 $x^2+y^2=(xy)^2$

 By Simon's Favorite Factoring Trick, we get $1=(1-x)(1+x)(1-y)(1+y)$

 Notice that other than for $x$ and/or $y$ equal to $-1$ and/or $1$ because for any $x$ or $y$ we can get some coordinates that work even if they're fractions which through $xy=k$ it can become an integer as $\dfrac{1}{2}\times2=1$

 This means we have to exclude $xy=k=1$ or $-1$. This is $\boxed{2}$ values that do not work.

 



\textbf{Solution 7 (Trig)}

The first equation is describing a circle with radius $k$ centered at $(0,0)$, so we can rewrite it in terms of a new variable $\theta$:

 
 \[(k\cos(\theta))^2+(k\sin(\theta))^2=k^2.\]

 A quick check shows that this identity holds since $\sin^2+\cos^2=1$. Now, with an alternate expression for $x$ and $y$, we can plug this into the second equation to get:

 
 \[(k\sin(\theta))(k\cos(\theta))=k\rightarrow k^2\sin(\theta)\cos(\theta)=k.\]

 Since the two equations does intersect when $k=0$ at $(x,y)=(0,0)$, we do not need to consider the $k=0$ case, so we can safely divide both sides by $k^2$ to get:

 
 \[\sin(\theta)\cos(\theta)=\frac{1}{k}.\]

 We can see that the intersections of these two equations will remain unchanged, only now that the solutions will be in terms of $k$ and $\theta$. Because the range of both $\sin$ and $\cos$ are limited to $[0,1]$, the range of $\sin(\theta)\cos(\theta)$ will not exceed $1$. However, the only way that the product could equal 1 is if both $\sin(\theta)$ and $\cos(\theta)$ to equal 1 at the same time, so since the maxima of the functions don't match up, this is an impossible case. Therefore, this equation will only have valid solutions when $k>1$, not at $k=1$. A quick sanity check with $\theta=\frac{\pi}{4}$ shows that $\frac{1}{k}=\frac{1}{2}$ and all values less than $\frac{1}{2}$ is indeed contained in the range.

 For all $k<0$, symmetry about the $y$-axis will give an identical argument with the equation $\sin(\theta)\cos(\theta)=-\frac{1}{k}$, with $k<-1$ producing valid solutions for our system. Then, the only two values of $k$ such that the system of equations won't intersect are $1$ and $-1$, a total of $\boxed{2}$ values.

 (Note: In the graph of $f(x)=\sin(x)\cos(x)$, the range of the function is actually even more restricted, being $\left[-\frac{1}{2},\frac{1}{2}\right]$. However, in a contest, observing that the value of $|\sin(x)\cos(x)|<1$ would suffice since the question is only asking for integer values of $k$.)

 


	\item (\textit{AMC 12 12A 2010}) Nondegenerate $\triangle ABC$ has integer side lengths, $\overline{BD}$ is an angle bisector, $AD = 3$, and $DC=8$. What is the smallest possible value of the perimeter?


 $\textbf{(A)}\ 30 \qquad \textbf{(B)}\ 33 \qquad \textbf{(C)}\ 35 \qquad \textbf{(D)}\ 36 \qquad \textbf{(E)}\ 37$

 \textbf{Answer: }B



\textbf{Solution}

By the \href{https://artofproblemsolving.com/wiki/index.php?title=Angle\_Bisector\_Theorem}{Angle Bisector Theorem}, we know that $\frac{AB}{BC} = \frac{3}{8}$. If we use the lowest possible integer values for $AB$ and $BC$ (the lengths of $AD$ and $DC$, respectively), then $AB + BC = AD + DC = AC$, contradicting the \href{https://artofproblemsolving.com/wiki/index.php?title=Triangle\_Inequality}{Triangle Inequality}. If we use the next lowest values ($AB = 6$ and $BC = 16$), the Triangle Inequality is satisfied. Therefore, our answer is $6 + 16 + 3 + 8 = \boxed{33}$, or choice $\textbf{(B)}$.

 



\textbf{Solution 2(Trick)}

We find that $\frac{AB}{BC}=\frac{3}{8}$ by the \href{https://artofproblemsolving.com/wiki/index.php?title=Angle\_Bisector\_Theorem}{Angle Bisector Theorem} so we let the lengths be $3n$ and $8n$, respectively where $n$ is a positive integer. Also since $AD=3$ and $BC=8$, we notice that the perimeter of the triangle is the sum of these, namely $3n+8n+3+8=11n+11.$ This can be factored into $11(n+1)$ and so the sum must be a multiple of $11$. The only answer choice which is a multiple of $11$ is $\boxed{\textbf{(B)} 33}$. ~mathboy282

 


	\item (\textit{AMC 12 12A 2010}) A coin is altered so that the probability that it lands on heads is less than $\frac{1}{2}$ and when the coin is flipped four times, the probability of an equal number of heads and tails is $\frac{1}{6}$. What is the probability that the coin lands on heads?


 $\textbf{(A)}\ \frac{\sqrt{15}-3}{6} \qquad \textbf{(B)}\ \frac{6-\sqrt{6\sqrt{6}+2}}{12} \qquad \textbf{(C)}\ \frac{\sqrt{2}-1}{2} \qquad \textbf{(D)}\ \frac{3-\sqrt{3}}{6} \qquad \textbf{(E)}\ \frac{\sqrt{3}-1}{2}$

 \textbf{Answer: }D



\textbf{Solution}

Let $x$ be the probability of flipping heads. It follows that the probability of flipping tails is $1-x$.

 The probability of flipping $2$ heads and $2$ tails is equal to the number of ways to flip it times the product of the probability of flipping each coin.

 
 \begin{align*}{4 \choose 2}x^2(1-x)^2 &= \frac{1}{6}\\ 6x^2(1-x)^2 &= \frac{1}{6}\\ x^2(1-x)^2 &= \frac{1}{36}\\ x(1-x) &= \pm\frac{1}{6}\end{align*}


 As for the desired probability $x$ both $x$ and $1-x$ are nonnegative, we only need to consider the positive root, hence

 
 \begin{align*}x(1-x) &= \frac{1}{6}\\ 6x^2-6x+1&=0\end{align*}


 Applying the quadratic formula we get that the roots of this equation are $\frac{3\pm\sqrt{3}}{6}$. As the probability of heads is less than $\frac{1}{2}$, we get that the answer is $\boxed{\textbf{(D)}\ \frac{3-\sqrt{3}}{6}}$.

 


	\item (\textit{AMC 12 12A 2010}) Bernardo randomly picks 3 distinct numbers from the set $\{1,2,3,...,7,8,9\}$ and arranges them in descending order to form a 3-digit number. Silvia randomly picks 3 distinct numbers from the set $\{1,2,3,...,6,7,8\}$ and also arranges them in descending order to form a 3-digit number. What is the probability that Bernardo's number is larger than Silvia's number?


 $\textbf{(A)}\ \frac{47}{72} \qquad \textbf{(B)}\ \frac{37}{56} \qquad \textbf{(C)}\ \frac{2}{3} \qquad \textbf{(D)}\ \frac{49}{72} \qquad \textbf{(E)}\ \frac{39}{56}$

 \textbf{Answer: }B



\textbf{Solution 1}

We can solve this by breaking the problem down into $2$ cases and adding up the probabilities.

 
Case $1$: Bernardo picks $9$.If Bernardo picks a $9$ then it is guaranteed that his number will be larger than Silvia's. The probability that he will pick a $9$ is $\frac{1 \cdot \binom{8}{2}}{\binom{9}{3}} = \frac{\frac{8\cdot7}{2}}{\frac{9\cdot8\cdot7}{3\cdot2\cdot1}}=\frac{1}{3}$.

 
Case $2$: Bernardo does not pick $9$.Since the chance of Bernardo picking $9$ is $\frac{1}{3}$, the probability of not picking $9$ is $\frac{2}{3}$.

 If Bernardo does not pick $9$, then he can pick any number from $1$ to $8$. Since Bernardo is picking from the same set of numbers as Silvia, the probability that Bernardo's number is larger is equal to the probability that Silvia's number is larger.

 Ignoring the $9$ for now, the probability that they will pick the same number is the number of ways to pick Bernardo's 3 numbers divided by the number of ways to pick any 3 numbers.

 We get this probability to be $\frac{3!}{8\cdot{7}\cdot{6}} = \frac{1}{56}$

 The probability of Bernardo's number being greater is
 \[\frac{1-\frac{1}{56}}{2} = \frac{55}{112}\]

 Factoring the fact that Bernardo could've picked a $9$ but didn't:

 
 \[\frac{2}{3}\cdot{\frac{55}{112}} = \frac{55}{168}\]

 Adding up the two cases we get $\frac{1}{3}+\frac{55}{168} = \boxed{\frac{37}{56}\ \textbf{(B)}}$



	\item (\textit{AMC 12 12A 2010}) Equiangular hexagon $ABCDEF$ has side lengths $AB=CD=EF=1$ and $BC=DE=FA=r$. The area of $\triangle ACE$ is $70\%$ of the area of the hexagon. What is the sum of all possible values of $r$?


 $\textbf{(A)}\ \frac{4\sqrt{3}}{3} \qquad \textbf{(B)} \frac{10}{3} \qquad \textbf{(C)}\ 4 \qquad \textbf{(D)}\ \frac{17}{4} \qquad \textbf{(E)}\ 6$

 \textbf{Answer: }E



\textbf{Solution 1}

It is clear that $\triangle ACE$ is an equilateral triangle. From the \href{https://artofproblemsolving.com/wiki/index.php?title=Law\_of\_Cosines}{Law of Cosines} on $\triangle ABC$, we get that $AC^2 = r^2+1^2-2r\cos{\frac{2\pi}{3}} = r^2+r+1$. Therefore, the area of $\triangle ACE$ is $\frac{\sqrt{3}}{4}(r^2+r+1)$ by area of an equilateral triangle.

 If we extend $BC$, $DE$ and $FA$ so that $FA$ and $BC$ meet at $X$, $BC$ and $DE$ meet at $Y$, and $DE$ and $FA$ meet at $Z$, we find that hexagon $ABCDEF$ is formed by taking equilateral triangle $XYZ$ of side length $r+2$ and removing three equilateral triangles, $ABX$, $CDY$ and $EFZ$, of side length $1$. The area of $ABCDEF$ is therefore

 $\frac{\sqrt{3}}{4}(r+2)^2-\frac{3\sqrt{3}}{4} = \frac{\sqrt{3}}{4}(r^2+4r+1)$.

 
Based on the initial conditions,

 
 \[\frac{\sqrt{3}}{4}(r^2+r+1) = \frac{7}{10}\left(\frac{\sqrt{3}}{4}\right)(r^2+4r+1)\]

 Simplifying this gives us $r^2-6r+1 = 0$. By \href{https://artofproblemsolving.com/wiki/index.php?title=Vieta\%27s\_Formulas}{Vieta's Formulas} we know that the sum of the possible value of $r$ is $\boxed{\textbf{(E)}\ 6}$.

 



\textbf{Solution 2}

Step 1: Use \href{https://artofproblemsolving.com/wiki/index.php?title=Law\_of\_Cosines}{Law of Cosines} in the same manner as the previous solution to get $AC=\sqrt{r^2+r+1}$. 

 Step 2: $\triangle{ABC}$~$\triangle{CDE}$~$\triangle{EFA}$ via SAS congruency. Using the formula $[ABC]=\frac{ab \sin C}{2}= \frac{r \sqrt{3}}{4}$. The area of the hexagon is equal to $[ACE] + 3[ABC]$. We are given that $70\%$ of this area is equal to $[ACE]$; solving for $AC$ in terms of $r$ gives $AC=\sqrt{7r}$.

 Step 3: $\sqrt{7r}=\sqrt{r^2+r+1} \implies r^2-6r+1=0$ and by \href{https://artofproblemsolving.com/wiki/index.php?title=Vieta\%27s\_Formulas}{Vieta's Formulas} , we get $\boxed{\textbf{E}}$.

 Note: To verify that the quadratic $r^2-6r+1$ has two positive roots, we can either solve for the roots directly or note that discriminant is positive, and there are no negative roots (because then $r^2, -6r, 1$ would all be positive).

 



\textbf{Solution 3}

Find the area of the triangle $ACE$ as how it was done in solution 1. Find the sum of the areas of the congruent  triangles $ABC, CDE, EFA$ as how it was done in solution 2. The sum of these areas is the area of the hexagon, hence the areas of the congruent  triangles $ABC, CDE, EFA$ is $30\%$ of the area of the hexagon. Hence $\frac{7}{3}$ times the latter is equal to the triangle $ACE$. Hence $\frac{7}{3}\cdot\frac{3\sqrt{3}}{4}r=\frac{\sqrt{3}}{4}(r^2+r+1)$. We can simplify this to $7r=r^2+r+1\implies r^2-6r+1=0$. By Vieta's, we get the sum of all possible values of $r$ is $-\frac{-6}{1}=6\text{ or } \boxed{\textbf{E}}$.-vsamc

 



\textbf{Solution 4 (no trig)}


                    \begin{figure}[H]
                        \centering
                        \includegraphics[width=0.5\textwidth]{Solution/AMC_12_12A_2010_17_Solution_Figure_1.png}
                    \end{figure}
                    Drawing not to scale.*To find r, we'll form an equation by finding the area of hexagon $ABCDEF$ and $\triangle{ACE}$ separately. First, connect a directly opposite diagonal such as $CF$. Because the hexagon is equiangular, edges $AF$ and $BC$ protrude from points $A$ and $B$ at the same angle but in opposite directions and equal length. This places points $C$ and $F$ at the same "height" relative to segment $AB$, so segment $CF$ is parallel to $AB$ and also $DE$ (opposite sides are parallel since you rotated through three $120^\circ$ angles between each other).

 Next, we'll calculate the area of isosceles trapezoids $ABCF$ and $CDEF$. Drop a perpendicular down from $B$ to $CF$, and call the intersection to $CF$, $P$. Because adjacent angles between parallel sides in a trapezoid sum to $180^\circ$ and $\angle ABC = 120^\circ$, $\angle BCP = 60^\circ$ meaning $\triangle{BPC}$ is a 30-60-90 right triangle. Now we get $BP = \frac{\sqrt{3}}{2}r$ and $CP = \frac{r}{2}$. $CF$ is equal to $AB + 2CP = r+1$. Hence, the area of trapezoid $ABCF$ is $\frac{1}{2} \cdot \frac{\sqrt{3}}{2}r \cdot (r+1+1) = \frac{\sqrt{3}}{4} r(r+2)$. Similarly, on trapezoid $CDEF$ we find $\triangle EQF$ is 30-60-90, and $EQ = \frac{\sqrt{3}}{2}, FQ = \frac{1}{2}$. Then, area of $CDEF$ is $\frac{\sqrt{3}}{4} (2r + 1)$. Thus, the area of the hexagon is the sum of the two areas, $\frac{\sqrt{3}}{4} (r^2 + 4r + 1)$.

 
Next, notice that $\triangle ACE$ is equilateral. And by Pythagorean theorem in $\triangle CQE$, side $CE = \sqrt{r^2 + r + 1}$. Applying area of equilateral triangle formula, area $\triangle ACE = \frac{\sqrt{3}}{4} (r^2 + r + 1)$. Using this area and the area given from 70\% of the area of the hexagon, we get the equation $\frac{7}{10} \cdot \frac{\sqrt{3}}{4} (r^2 + 4r + 1) = \frac{\sqrt{3}}{4} (r^2 + r + 1)$. Simplifying, we get the quadratic $r^2-6r+1=0$, which by Vieta's yields the answer $6, \boxed{\textbf{E}}$.

 \textbf{Proof Triangle ACE is Equilateral.} We know $\triangle{ABC}$, $\triangle{CDE}$, and $\triangle{EFA}$ are congruent by SAS, so the side opposite the 120 degree angle is also the same (since the triangles are congruent). Thus $\triangle{ACE}$ is equilateral.Q.E.D.~mathboy282

 


	\item (\textit{AMC 12 12A 2010}) A 16-step path is to go from $(-4,-4)$ to $(4,4)$ with each step increasing either the $x$-coordinate or the $y$-coordinate by 1. How many such paths stay outside or on the boundary of the square $-2 \le x \le 2$, $-2 \le y \le 2$ at each step?


 $\textbf{(A)}\ 92 \qquad \textbf{(B)}\ 144 \qquad \textbf{(C)}\ 1568 \qquad \textbf{(D)}\ 1698 \qquad \textbf{(E)}\ 12,800$

 \textbf{Answer: }D



\textbf{Solution 1}

Each path must go through either the second or the fourth quadrant.Each path that goes through the second quadrant must pass through exactly one of the points $(-4,4)$, $(-3,3)$, and $(-2,2)$.

 There is $1$ path of the first kind, ${8\choose 1}^2=64$ paths of the second kind, and ${8\choose 2}^2=28^2=784$ paths of the third type. Each path that goes through the fourth quadrant must pass through exactly one of the points $(4,-4)$, $(3,-3)$, and $(2,-2)$.Again, there is $1$ path of the first kind, ${8\choose 1}^2=64$ paths of the second kind, and ${8\choose 2}^2=28^2=784$ paths of the third type. 

 Hence the total number of paths is $2(1+64+784) = \boxed{1698}$.

 



\textbf{Solution 2}

We will use the concept of complementary counting. If you go in the square you have to go out by these labeled points(click the link) \href{https://imgur.com/VysX4P0}{https://imgur.com/VysX4P0} and go out through the borders because if you didn't, you would touch another point in one of those points in the set of points in the link(Call it $S$). There is symmetry about $y=x$, so we only have to consider $(1,-1), (1,0)$, and $(1,1)$. $(1,1)$ can go on the boundary in two ways, so we can only consider one case and multiply it by two. For $(1,0)$ and $(1,-1)$ we can just multiply by two. So we count paths from $(-4,4)$ to each of these points, and then multiply that by the number of ways to get from the point one unit right of that to $(4,4)$, and all in all, we get the answer is $\dbinom{16}{8}-2\left[\dbinom{8}{3}\dbinom{7}{2}+\dbinom{9}{4}\dbinom{6}{2}+\dbinom{10}{5}\dbinom{5}{2}\right]=1698$, which is answer choice $\textbf{(D)}$-vsamcNote: Sorry if this was rushed.

 \textbf{Further Explanation(for Solution 1)} As stated in the solution, there are $6$ points along the line $y=-x$ that constitute a sort of "boundary". Once the ant reaches one of these $6$ points, it is exactly halfway to $(4, 4)$. Also notice that the ant will only cross one of the $6$ points during any one of its paths. Therefore we can divide the problem into $3$ cases, focusing on $1$ quadrant; then multiplying the sum by $2$ to get the total (because there is symmetry).

 
For the sake of this explanation, we will focus on the fourth quadrant (it really doesn't matter which quadrant because, again the layout is symmetrical) The three cases are when the ant crosses $(4, -4), (3, -3),$ and $(2, -2)$.

 
For each of the cases, notice that the path the ant takes can be expressed as a sequence of steps, such as:

 
right, right, up, right,..., etc.

 
Also notice that there are always $8$ steps per sequence (if there were more or less steps, the ant would be breaking the conditions given in the problem). This means we can figure out the number of ways to get to a point based on the particular sequence of steps that denote each path. For example, there is only one way for the ant to pass through $(4, -4);$ it MUST keep traveling right for all $8$ steps. This seems fairly obvious; however, notice that this is equivalent to

 
$\binom{8}{0}$

 
Now we consider the number of ways to get from $(4, -4)$ to $(4, 4)$. by symmetry, there is only $1$ such way. So the number of paths containing $(4, -4)$ is $1^2,$ or $1$.

 
Moving on to the next case, we see that the ant MUST travel right exactly $7$ times and up exactly once. So each sequence of this type will have $7$ "right"s and $1$ "up". So, the total number of paths that go through $(3, -3)$ is equivalent to the number of ways to arrange $1$ "up" into $8$ spots. This is

 
$\binom{8}{1} = 8$

 
Similarly to the first case, we square this value to account for the second half of the journey: $8^2 = 64$.

 
Finally, for the third case (ant passes through $(2, -2)$) the ant must travel right exactly $6$ times and up exactly $2$ times. This is equivalent to the number of ways to arrange $2$ "ups" in a sequence of $8$ movements, or

 
$\binom{8}{2} = 28$

 
Again, we square $28$: $28^2 = 784$. Adding up all of these cases we get

 
$1+64+784 = 849$

 
paths through the fourth quadrant. Doubling this number to account for the paths through the second quadrant, we have

 
$849*2=1698 \Rightarrow \boxed{\text{D}}$.

 
For all the notation geeks out there, the solution can be expressed as such:

 
$2\sum_{k=0}^2 \binom{8}{k}^2$




\textbf{Solution 3}

We can draw an $8$ by $8$ square with a hollow $4$ by $4$ center. The ways to get a to a point or corner of a coordinate point is equal to the ways of getting to the point to the left of the desired point and the bottom of the desired point, since we can only move right and up. Using basic addition and symmetry (along the $y=x$) to speed things up, in the end, you sum up $849$ and $849,$ giving us our answer of $\boxed{1698}.$

                    \begin{figure}[H]
                        \centering
                        \includegraphics[width=0.5\textwidth]{Solution/AMC_12_12A_2010_18_Solution_Figure_1.png}
                    \end{figure}
                    




\textbf{Solution 4}

Let $N(A,C,B)$ denote the number of paths from $A$ to $C$ then to $B$,  $N(A,D,B)$ denote the number of paths from $A$ to $D$ then to $B$, $N(A,E,B)$ denote the number of paths from $A$ to $E$ then to $B$, $N(A,C,D, B)$ denote the number of paths from $A$ to $C$ to $D$ then to $B$, $N(A,C,D,E,B)$ denote the number of paths from $A$ to $C$  to $D$ to $E$ then to $B$, also notice we have $N(A,C,E,B)=N(A,C,D,E,B)$.
                    \begin{figure}[H]
                        \centering
                        \includegraphics[width=0.5\textwidth]{Solution/AMC_12_12A_2010_18_Solution_Figure_2.png}
                    \end{figure}
                    then the answer = $N(A,C,B)+N(A,D,B)+N(A,E,B)-N(A,C,D,B)-N(A,C,E,B)-N(A,D,E,B)+N(A,C,D,E,B)=N(A,C,B)+N(A,D,B)+N(A,E,B)-N(A,C,D,B)-N(A,D,E,B)$, $N(A,C,B)=N(A,C)\cdot N(C,B)={8\choose 2}\cdot {8\choose 2}$, $N(A, D, B)=N(A,D)\cdot N(D, B)={9\choose 2}\cdot {7\choose 1}$, $N(A,E, B)={10\choose 2}\cdot 1$, $N(A, C, D, B)={8\choose 2}\cdot {7\choose 1}, N(A, D, E, B)={9\choose 2}\cdot1,$ since if the path passes quadrant II, it will pass either $C$ or $D$ or $E$. Consider the paths pass the quadrant IV together, we have ans=$2\times({8\choose 2}\cdot {8\choose 2}+{9\choose 2}\cdot {7\choose 1}+{10\choose 2}\cdot 1-{8\choose 2}\cdot {7\choose 1}-{9\choose 2}\cdot1)=2\times 849=1698$,$\boxed{(D)}$



	\item (\textit{AMC 12 12A 2010}) Each of 2010 boxes in a line contains a single red marble, and for $1 \le k \le 2010$, the box in the $k\text{th}$ position also contains $k$ white marbles. Isabella begins at the first box and successively draws a single marble at random from each box, in order. She stops when she first draws a red marble. Let $P(n)$ be the probability that Isabella stops after drawing exactly $n$ marbles. What is the smallest value of $n$ for which $P(n) < \frac{1}{2010}$?


 $\textbf{(A)}\ 45 \qquad \textbf{(B)}\ 63 \qquad \textbf{(C)}\ 64 \qquad \textbf{(D)}\ 201 \qquad \textbf{(E)}\ 1005$

 \textbf{Answer: }A



\textbf{Solution 1}

The probability of drawing a white marble from box $k$ is $\frac{k}{k + 1}$, and the probability of drawing a red marble from box $k$ is $\frac{1}{k+1}$.

 To stop after drawing $n$ marbles, we must draw a white marble from boxes $1, 2, \ldots, n-1,$ and draw a red marble from box $n.$ Thus, 
 \[P(n) = \left(\frac{1}{2} \cdot \frac{2}{3} \cdot \frac{3}{4} \cdots \frac {n - 1}{n}\right) \cdot \frac{1}{n +1} = \frac{1}{n (n + 1)}.\]

 So, we must have $\frac{1}{n(n + 1)} < \frac{1}{2010}$ or $n(n+1) > 2010.$

 Since $n(n+1)$ increases as $n$ increases, we can simply test values of $n$; after some trial and error, we get that the minimum value of $n$ is $\boxed{\textbf{(A) }45}$, since $45(46) = 2070$ but $44(45) = 1980.$




\textbf{Solution 2(cheap)}

Do the same thing as Solution 1, but when we get to $n(n+1)>2010$ just test all the answer choices in ascending order(from A to E), and stop when one of the answer choices is greater than $2010$. We get $45(46)=2070$, which is greater than $2010$, so we are done. The answer is $\textbf{(A)}$

 -vsamc

 



\textbf{Solution 3(very cheap)}

Do the same thing as Solution 1, but when we get to $n(n+1)>2010$, remember that in our current era, the square root of the year is in the 40s. $n(n+1)$ is close to $n^2$, so only $\textbf{(A)}$ can be correct.

 


	\item (\textit{AMC 12 12A 2010}) Arithmetic sequences $\left(a_n\right)$ and $\left(b_n\right)$ have integer terms with $a_1=b_1=1<a_2 \le b_2$ and $a_n b_n = 2010$ for some $n$. What is the largest possible value of $n$?


 $\textbf{(A)}\ 2 \qquad \textbf{(B)}\ 3 \qquad \textbf{(C)}\ 8 \qquad \textbf{(D)}\ 288 \qquad \textbf{(E)}\ 2009$

 \textbf{Answer: }C



\textbf{Solution 1}

Since $\left(a_n\right)$ and $\left(b_n\right)$ have integer terms with $a_1=b_1=1$, we can write the terms of each sequence as

 
 \begin{align*}&\left(a_n\right) \Rightarrow \{1, x+1, 2x+1, 3x+1, ...\}\\ &\left(b_n\right) \Rightarrow \{1, y+1, 2y+1, 3y+1, ...\}\end{align*}


 where $x$ and $y$ ($x\leq y$) are the common differences of each, respectively.

 
Since

 
 \begin{align*}a_n &= (n-1)x+1\\ b_n &= (n-1)y+1\end{align*}


 it is easy to see that

 $a_n \equiv b_n \equiv 1 \mod{(n-1)}$.

 
Hence, we have to find the largest $n$ such that $\frac{a_n-1}{n-1}$ and $\frac{b_n-1}{n-1}$ are both integers; equivalently, we want to maximize $\gcd(a_n-1, b_n-1)$.

 
The prime factorization of $2010$ is $2\cdot{3}\cdot{5}\cdot{67}$. We list out all the possible pairs that have a product of $2010$, noting that these are the possible values of $(a_n, b_n)$ and we need $a_n \leq b_n$:

 
 \[(2,1005), (3, 670), (5,402), (6,335), (10,201),(15,134),(30,67)\]

 and soon find that the largest $n-1$ value is $7$ for the pair $(15, 134)$, and so the largest $n$ value is $\boxed{8\ \textbf{(C)}}$.

 



\textbf{Solution 2}

As above, let $a_n=(n-1)x+1$ and $b_n=(n-1)y+1$ for some $1\leq x\leq y$.

 Now we get $2010 = a_n b_n = (n-1)^2xy + (n-1)(x+y) + 1$, hence $2009 = (n-1)( (n-1)xy + x + y )$. Therefore $n-1$ divides $2009 = 7^2 \cdot 41$. And as the second term is greater than the first one, we only have to consider the options $n-1\in\{1,7,41\}$.

 For $n=42$ we easily see that for $x=y=1$ the right side is less than $49$ and for any other $(x,y)$ it is way too large.

 For $n=8$ we are looking for $(x,y)$ such that $7xy + x + y = 2009/7 = 7\cdot 41$. Note that $x+y$ must be divisible by $7$. We can start looking for the solution by trying the possible values for $x+y$, and we easily discover that for $x+y=21$ we get $xy + 3 = 41$, which has a suitable solution $(x,y)=(2,19)$.

 Hence $n=8$ is the largest possible $n$. (There is no need to check $n=2$ anymore.)

 



\textbf{Solution 3 (using answer choices)}

Consider $n=288$, which would imply $b_{288}\ge a_{288}\ge 288$. However then $a_n b_n\ge 288^2>2010$, so we just need to show that $n=8$ is achievable. This is true when $a_n=1+2n$ and $b_n=1+19n$, giving $a_8 b_8=(15)(134)=2010$. Hence the answer is $\boxed{\textbf{(C)}\ 8}$.

 \textbf{Alternative Thinking} Since

 
$a_n\times b_n = 2010,$

 
and

 
$a_n \le b_n$,

 


 it follows that

 
$a_n \le \sqrt{2010} \Rightarrow a_n \le 44$.

 
But $a_n$ and $b_n$ are also integers, so $a_n$ must be a factor of $2010$ smaller than $44$. Notice that $2010 = 2*3*5*67$. Therefore $a_n = 2, 3, 5, 6, 112, 15,$ or $30$ and $b_n = 1005, 670, 402, 335, 201, 134,$ or $67$; respectively.

 
Notice that the term $a_m$ is equivalent to the first term $a_1 = 1$ plus $(m-1)$ times the common difference for that particular arithmetic sequence. Let the common difference of $(a_n)$ be $k$ and the common difference of $(b_n)$ be $i$ (not $\sqrt{-1}$). Then

 
$a_n$ (the $n$th term, not the sequence itself) $=1 + k(n-1)$

 
and

 
$b_n = 1 + i(n-1)$

 
Subtracting one from all the possible values listed above for $a_n$ and $b_n$, we get

 
$k(n-1) = 1, 2, 4, 5, 9, 14, 29$

 
and

 
$i(n-1) = 1004, 669, 401, 334, 200, 133, 66$

 
In order to maximize $n$, we must maximize $n-1$. Therefore $k$ and $i$ are \href{https://artofproblemsolving.com/wiki/index.php?title=Coprime}{coprime} and $n-1$ is the \href{https://artofproblemsolving.com/wiki/index.php?title=Greatest\_common\_factor}{GCF} of any corresponding pair. Inspecting all of the pairs, we see that the \href{https://artofproblemsolving.com/wiki/index.php?title=Greatest\_common\_factor}{GCF} is always $1$ except for the pair $(14, 133),$ which has a GCF of $7$. Therefore the maximum value of $n$ is $8 \Rightarrow \boxed{\text{C}}$.

 


	\item (\textit{AMC 12 12A 2010}) The graph of $y=x^6-10x^5+29x^4-4x^3+ax^2$ lies above the line $y=bx+c$ except at three values of $x$, where the graph and the line intersect. What is the largest of these values?


 $\textbf{(A)}\ 4 \qquad \textbf{(B)}\ 5 \qquad \textbf{(C)}\ 6 \qquad \textbf{(D)}\ 7 \qquad \textbf{(E)}\ 8$

 \textbf{Answer: }A



\textbf{Solution 1}

The $x$ values in which $y=x^6-10x^5+29x^4-4x^3+ax^2$ intersect at $y=bx+c$ are the same as the zeros of $y=x^6-10x^5+29x^4-4x^3+ax^2-bx-c$.

 Since there are $3$ zeros and the function is never negative, all $3$ zeros must be double roots because the function's degree is $6$.

 Suppose we let $p$, $q$, and $r$ be the roots of this function, and let $x^3-ux^2+vx-w$ be the cubic polynomial with roots $p$, $q$, and $r$.

 
 \begin{align*}(x-p)(x-q)(x-r) &= x^3-ux^2+vx-w\\ (x-p)^2(x-q)^2(x-r)^2 &= x^6-10x^5+29x^4-4x^3+ax^2-bx-c = 0\\ \sqrt{x^6-10x^5+29x^4-4x^3+ax^2-bx-c} &= x^3-ux^2+vx-w = 0\end{align*}


 In order to find $\sqrt{x^6-10x^5+29x^4-4x^3+ax^2-bx-c}$ we must first expand out the terms of $(x^3-ux^2+vx-w)^2$.

 
 \[(x^3-ux^2+vx-w)^2\]
 \[= x^6-2ux^5+(u^2+2v)x^4-(2uv+2w)x^3+(2uw+v^2)x^2-2vwx+w^2\]

 [Quick note: Since we don't know $a$, $b$, and $c$, we really don't even need the last 3 terms of the expansion.]

 
 \begin{align*}&2u = 10\\ u^2+2v &= 29\\ 2uv+2w &= 4\\ u &= 5\\ v &= 2\\ w &= -8\\ &\sqrt{x^6-10x^5+29x^4-4x^3+ax^2-bx-c} = x^3-5x^2+2x+8\end{align*}


 All that's left is to find the largest root of $x^3-5x^2+2x+8$.

 
 \begin{align*}&x^3-5x^2+2x+8 = (x-4)(x-2)(x+1)\\ &\boxed{\textbf{(A)}\ 4}\end{align*}





\textbf{Solution 2}

The $x$ values in which $y=x^6-10x^5+29x^4-4x^3+ax^2$ intersect at $y=bx+c$ are the same as the zeros of $y=x^6-10x^5+29x^4-4x^3+ax^2-bx-c$.We also know that this graph has 3 places tangent to the x-axis, which means that each root has to have a multiplicity of 2.Let the function be $(x-p)^2(x-q)^2(x-r)^2$. 

 Applying Vieta's formulas, we get $2p+2q+2r = 10$  or  $p+q+r = 5$.Applying it again, we get, after simplification, $p^2+q^2+r^2+4pq+4pr+4qr = 29$. 

 Notice that squaring the first equation yields $p^2+q^2+r^2+2pq+2qr+2pr= 25$, which is similar to the second equation.

 Subtracting this from the second equation, we get $2pq+2pr+2qr = 4$. Now that we have the $pq+pr+qr$ term, we can manpulate the equations to yield the sum of squares. $2(p^2+q^2+r^2+2pq+2qr+2pr)-2pq-2pr-2qr= 25*2-4$ or  $2p^2+2q^2+2r^2+2pq+2qr+2pr = 46$. We finally reach $(p+q)^2+(q+r)^2+(p+r)^2 = 46$. 

 Since the answer choices are integers, we can guess and check squares to get $\{(p+q)^2, (q+r)^2, (p+r)^2\} = \{1, 9, 36\}$ in some order. We can check that this works by adding then and seeing $2p+2q+2r = 10$. We just need to take the lowest value in the set, square root it, and subtract the resulting value from 5 to get $\boxed{\textbf{(A)}\ 4}$.

 Note: One could also multiply $2pq+2pr+2qr = 4$ by 2 and subtract from $p^2+q^2+r^2+4pq+4pr+4qr = 29$ to obtain 
 \[p^2+q^2+r^2=21\] The ordered triple {16,4,1} sums to 21, and the answer choices are all positive integers, therefore the answer is 4. 

 \textbf{Alternative method:} After reaching $p+q+r = 5$ and $pq + qr + rp = 2$, we can algebraically derive $pqr$.

 Applying Vieta's formulas on the $x^3$ term yields $2p^2q+2pq^2+2q^2r+2qr^2+2r^2p+2rp^2+8pqr = 4$.

 Notice that $(p+q+r)(pq+qr+rp) = p^2q+pq^2+q^2r+qr^2+r^2p+rp^2+3pqr$, so 
 \[2p^2q+2pq^2+2q^2r+2qr^2+2r^2p+2rp^2+6pqr = 2(p+q+r)(pq+qr+rp) = 20.\]

 Subtracting this from $2p^2q+2pq^2+2q^2r+2qr^2+2r^2p+2rp^2+8pqr = 4$ yields $2pqr = -16$, so $pqr = -8$, which means that $p$, $q$, and $r$ are the roots of the cubic $x^3 - 5x^2 + 2x + 8$, and it is not hard to find that these roots are $-1$, $2$, and $4$. The largest of these values is $\boxed{\textbf{(A)}\ 4}$.

 Extra note: One can also use the identity
 \[(p+q)(q+r)(p+r)=2pqr+\sum_{sym}p^2q\]to compute $pqr$. From $p^2q+pq^2+q^2r+qr^2+r^2p+rp^2+4pqr = 2$, use our identity to get $2pqr+(p+q)(q+r)(p+r)=2$. Then use $p+q+r=5$ to rewrite as $2pqr+(5-p)(5-q)(5-r)=2$. Expanding and using $pq+qr+pr=2$ as well gives the result $pqr=-8$.

 



\textbf{Solution 3}

The $x$-values where $y = x^6 - 10x^5 + 29x^4 - 4x^3 + ax^2$ intersects $y = bx + c$ are the zeros of

 $g(x) = x^6 - 10x^5 + 29x^4 - 4x^3 + ax^2 - bx - c$.

 We are told the graph lies above the line except at exactly three values of $x$. Thus $g(x) \ge 0$ for all $x$, and it has exactly three real roots.

 Since $g(x)$ has degree 6 and is always nonnegative, each real root must have even multiplicity. Because there are exactly three real roots and the total multiplicity is 6, each root has multiplicity 2.

 So we may write

 $g(x) = (x-p)^2(x-q)^2(x-r)^2$.

 Since the leading coefficient is 1, this equals

 $(x^3 + ux^2 + vx + w)^2$.

 Expanding,

 $(x^3 + ux^2 + vx + w)^2 = x^6 + 2u x^5 + (u^2 + 2v)x^4 + (2uv + 2w)x^3 + \cdots$

 Matching coefficients with$x^6 - 10x^5 + 29x^4 - 4x^3 + \cdots$ gives:

 $2u = -10 \Rightarrow u = -5$.

 $u^2 + 2v = 29 \Rightarrow 25 + 2v = 29 \Rightarrow v = 2$.

 $2uv + 2w = -4 \Rightarrow 2(-5)(2) + 2w = -4 \Rightarrow -20 + 2w = -4 \Rightarrow 2w = 16 \Rightarrow w = 8$.

 Thus,

 $g(x) = (x^3 - 5x^2 + 2x + 8)^2$.

 Now factor the cubic. Testing rational roots, we find $x=4$ works:

 $4^3 - 5(4^2) + 2(4) + 8 = 64 - 80 + 8 + 8 = 0$.

 Dividing gives

 $x^3 - 5x^2 + 2x + 8 = (x-4)(x^2 - x - 2)$.

 Factor the quadratic:

 $x^2 - x - 2 = (x-2)(x+1)$.

 So

 $g(x) = (x-4)^2 (x-2)^2 (x+1)^2$.

 The intersection points are $x = -1, 2, 4$.

 The largest is $\boxed{4}$.

 



\textbf{Solution 4}

First, $y=x^6-10x^5+29x^4-4x^3+ax^2-bx-c = 0$ has exactly $3$ roots. Therefore, $y = (kx^3+lx^2+mx+n)^2 = 0$.

 So, $k^2x^6+2klx^5+(2km+l^2)x^4+2(kn+lm)x^3+ax^2-bx-c = 0$

 By matching the coefficients of the first $4$ terms, we have $k^2 = 1, 2kl = -10, 2km+l^2 = 29, 2kn+2lm = -4$

 Solving the equations above, we have $2$ sets of solutions; first set of which is $k = 1, l = -5, m = 2, n = 8$. Second set of which is $k = -1, l = 5, m = -2, n = -8$. After squaring both sets, they are the same i.e. $x^3-5x^2+2x+8 = 0$.

 This is equal to $(x-4)(x-2)(x+1) = 0$. The largest root is $\boxed {\textbf{(A) 4}}$




\textbf{Solution 5}

The $x$ values in the coordinates of which these $2$ polynomials intersect is the same as the zeroes of
 \[P(x)=x^6-10x^5+29x^4-4x^3+ax^2-bx-c.\]Since there are only $3$ intersections, they must all be double roots. So, we let $r,s,$ and $t$ be the roots of $P(x)$. We let,
 \[(x-r)(x-s)(x-t)=x^3-ux^2+vx-w\]
 \[\Longrightarrow [(x-r)(x-s)(x-t)]^2=x^6-10x^5+29x^4-4x^3+ax^2-bx-c\]
 \[\Longrightarrow \sqrt{x^6-10x^5+29x^4-4x^3+ax^2-bx-c} = x^3-ux^2+vx-w\]
 \[\Longrightarrow x^6-10x^5+29x^4-4x^3+ax^2-bx-c=(x^3-ux^2+vx-w)^2\]
 \[\Longrightarrow x^6-10x^5+29x^4-4x^3+ax^2-bx-c=x^6-(2u)x^5+(u^2+2v)x^4-(2w+2uv)x^3+...\]We don't really care about the other coefficients because they aren't known. So, comparing coefficients, we get $(u,v,w)=(5,2,-8)$. So, we have
 \[(x-r)(x-s)(x-t)=x^3-5x^2+2x+8.\]Now, we can use RRT to finish from here. Factoring the polynomial gives
 \[x^3-5x^2+2x+8=(x+1)(x-2)(x-4),\]so, our largets root is $\boxed{\textbf{(A)}\ 4}.$

 -jb2015007

 


	\item (\textit{AMC 12 12A 2010}) What is the minimum value of $\left|x-1\right| + \left|2x-1\right| + \left|3x-1\right| + \cdots + \left|119x - 1 \right|$?


 $\textbf{(A)}\ 49 \qquad \textbf{(B)}\ 50 \qquad \textbf{(C)}\ 51 \qquad \textbf{(D)}\ 52 \qquad \textbf{(E)}\ 53$

 \textbf{Answer: }A



\textbf{Solution 1}

If we graph each term separately, we will notice that all of the zeros occur at $\frac{1}{m}$, where $m$ is any integer from $1$ to $119$, inclusive: $|mx-1|=0\implies mx=1\implies x=\frac{1}{m}$.

 The minimum value of $f(x)$ occurs where the absolute value of the sum of the slopes is at a minimum $\ge 0$, since it is easy to see that the value will be increasing on either side. That means the minimum must happen at some $\frac{1}{m}$.

 The sum of the slopes at $x = \frac{1}{m}$ is

 
 \begin{align*}&\sum_{i=m+1}^{119}i - \sum_{i=1}^{m}i\\ &=\sum_{i=1}^{119}i - 2\sum_{i=1}^{m}i\\ &=-m^2-m+7140\end{align*}


 Now we want to minimize $-m^2-m+7140$. The zeros occur at $-85$ and $84$, which means the slope is $0$ where $m = 84, 85$.

 We can now verify that both $x=\frac{1}{84}$ and $x=\frac{1}{85}$ yield $\boxed{49\ \textbf{(A)}}$.

 You can also think of the slopes playing 'tug of war', where the slope of each absolute function upon passing its $x$-intercept is negated, positively tugging on the remaining negative slopes.

 The sum of the slopes is $1+2+3+4\dots 119=\sum_{m=1}^{119}m=\frac{119\cdot 120}{2}=60\cdot 119=7140$

 So we need to find the least integer $a$ such that $1+2+3+\dots a=\sum_{n=1}^an=\frac{a(a+1)}{2}\ge \frac{7140}{2}=3570:$

 \[a(a+1)\ge 7140\implies a^2+a-7140\ge 0\rightarrow a=84\text{ exactly!}\]This "exactly" means that the slope is ZERO between the whole interval $x\in\left(\frac{1}{85},\frac{1}{84}\right)$. We can explicitly evaluate both to check that they are both equal to the desired minimum value of $f(x)$:

 
 \[\frac{84+83+\dots+2+1+1+2+\dots+33+34}{85}=\frac{84(85)/2+34(35)/2}{85}=\frac{85(14+84)/2}{85}=49\]

 \[\frac{83+82+\dots+2+1+1+2+\dots+34+35}{84}=\frac{83(84)/2+35(36)/2}{84}=\frac{84(15+83)/2}{84}=49\]

 Thus the minimum value of $f(x)$ is $49$.

 



\textbf{Solution 2}

Rewrite the given expression as follows:
 \[1|x-1| + 2\left|x-\frac 12\right| + \cdots + 119\left|x-\frac 1{119}\right|\]Imagine the real line. For each $n\in\{1,\dots,119\}$ imagine that there are $n$ boys standing at the coordinate $\frac 1n$. We now need to place a donut on the real line in such a way that the sum of its distances from all the boys is minimal, and we need to compute this sum.

 Note that there are $B=1+2+\cdots+119 = 119\cdot 60=7140$ boys in total. Let's label them from 1 (the only boy placed at $1$) to $B$ (the last boy placed at $\frac 1{119}$.

 Clearly, the minimum sum is achieved if the donut's coordinate is the median of the boys' coordinates. To prove this, place the donut at the median coordinate. If you now move it in any direction by any amount $d$, there will be $B/2$ boys such that it moves $d$ away from this boy. For each of the remaining boys, it moves at most $d$ closer, hence the total sum of distances does not decrease.

 Hence the optimal solution is to place the donut at the median coordinate. Or, more precisely, as $B$ is even, we can place it anywhere on the segment formed by boy $B/2$ and boy $(B/2)+1$: by extending the previous argument, anywhere on this segment the sum of distances is the same.

 By trial and error, or by solving the quadratic equation $z(z+1)/2 = 7140/2$ we get that boy number $B/2$ is the last boy placed at $\frac 1{84}$ and the next boy is the one placed at $\frac 1{85}$. Hence the given expression is minimized for any $x\in\left[ \frac 1{85}, \frac 1{84} \right]$.

 \textbf{Common part of both solutions} To find the minimum, we want to balance the expression so that it is neither top nor bottom heavy. $\frac{119(120)}{2(2)}=\frac{7140}{2}=3570=\frac{84(85)}{2}=\frac{119(120)}{2}-\frac{84(85)}{2}$.

 Now that we know that the sum of the first 84 $x$'s is equivalent to the sum of $x$'s 85 to 119, we can plug either $\frac{1}{84}$ or $\frac{1}{85}$ to find the minimum.

 Note that the terms $x-1$ to $83x-1$ are negative, and the terms $85x-1$ to $119x-1$ are positive. Hence we get:
 \begin{align*} & |x-1| + |2x-1| + \cdots + |83x-1| \\ =~ & (1-x) + (1-2x) + \cdots + (1-83x) \\ =~ & 83 - x(1+2+\cdots+83) \\ =~ & 83 - \frac 1{84} \cdot \frac{83\cdot 84}2 \\ =~ & 83 - \frac{83}2 \\ =~ & \frac{83}2 \end{align*}
and
 \begin{align*} & |85x-1| + |86x-1| + \cdots + |119x-1| \\ =~ & (85x-1) + (86x-1) + \cdots + (119x-1) \\ =~ & x(85+86+\cdots+119) - (119-84) \\ =~ & \frac 1{84} \cdot \frac{84\cdot 85}2 - 35 \\ =~ & \frac{85}2 - 35 \\ =~ & \frac{15}2 \end{align*}
Hence the total sum of distances is $\frac{83}2 + \frac{15}2 = 49$.

 



\textbf{Solution 3}

By the triangle inequality,$|x-1|+|2x-1|+|3x-1|+\cdots + |119x-1| \geq |(x-1)+(2x-1)+\cdots+(119x-1)|.$However, we may change signs of some of these terms to cancel out the $x$'s.Since the minimum exists, we want all the $x$s to cancel out. Thus, we want to find some $n$ such that 
 \[1+2+3+...+n=(n+1)+(n+2)+(n+3)+...+119\] 
 \[\frac{n(n+1)}{2}=\frac{119\cdot120}{2}-\frac{n(n+1)}{2}\] 
 \[n^2+n-7140=0\] 
 \[n=84\]

 Then, $x=\frac{1}{n}= \frac{1}{84}$. The answer(expression's value) is then $84*1+(119-85+1)*(-1)$ , which becomes $84-35=\boxed{49}$.

 



\textbf{Solution 4}

If $nx\geq1$ for some value $n$ between $1$ and $119$, we have that the $nx, (n+1)x, ... , 119x$ terms will be positive and $(n-1)x, (n-2)x, ... , 1x$ terms will be negative. 

 
Therefore, the expression given in the problem will be $nx-1 + (n+1)x-1 + ... + 119x-1 - ((n-1)x-1 + (n-2)x-1 + ... + x-1)$ which simplifies to $(7140-n(n-1))x + (2n-121).$

 The minimum value of x that achieves this state is when $nx=1,$ so we have $x=\frac{1}{n}.$ 

 Applying $x=\frac{1}{n}$ to the expression above, we have $\frac{7140}{n} + n - 120.$ 

 
To finish finding the minimizing $n$ of this expression, we either use AM-GM on $\frac{7140}{n} + n$ or derivatives.

 Either way, we find that the value of $n$ that minimizes is $84.49$ which should mean $n=85$. Plugging this value in, we get $\boxed{49}$ as our minimum value.

 


 \textbf{Comments} There is no motivation behind wanting all the $x$s to cancel out other than assumption that that would lead to the minimum.  There are two possibilities for the minima of a function such as this. One occurs when there is a flat region on the graph of this function as the minimum, which is when you can use the above solution, but the other one occurs when the minimum is achieved at a "kink", when the graph's slope switches from positive to negative at a single point. For example, if there were $7$ terms, since $\frac{7*8}{2}$ is not of the form $n(n+1)$, you cannot cancel all the $x$s. In fact, the unique global minimum would be achieved at $x=\frac{1}{5}$, where the slope jumps from negative to positive on a graph. Thus the assumption that the minimum is achieved when all the $x$s cancel is not valid.

 



\textbf{Solution 5}

We can take the derivative and set the equation equal to zero (stay with me now)

 The derivative of $|x| = s(x)$ (the sign (sgn) function), where $s(-x) = -1, s(0) = 0, s(x) = 1, \forall x \in \mathbf(R)^+$.

 Taking the derivative (with the chain rule) and setting the expression equal to zero yields

 $s(x-1) + 2s(2x-2) + \cdots + 119s(119x-1) = 0$.

 (I know it looks like s is a variable but I can't get the sgn function to render so if someone could fix that that would be great thank you)

 
Now we have two cases: either one of the terms is zero or they're all nonzero. Lets just assume one of them is zero because one of the other solutions said that if you moved left or right from a zero then you'd be getting bigger and it doesn't really matter anyway. 

 Now we can see that the expression is gonna look something like (a bunch of negative terms) + 0 + (a bunch of positive terms). 

 The negative terms are gonna add up to $-(1 + 2 + 3 + \cdots + n-1) = -\frac{n(n-1)}{2}$. The positive terms are gonna add up to $(n+1) + (n+2) + \cdots + 119 = \frac{119 \cdot 120}{2} - \frac{n(n-1)}{2}$. (it really should be n(n+1) but n=0 so it's all good)

 Ok now we have $119 \cdot 60 - 2\frac{n(n-1)}{2} = 0 \implies 7140 = n(n+1) \implies n = 84 or n = -85$. If we plug in $x=\frac{1}{84}$ we get $\boxed{49}$ and I think if we plug in $x = -\frac{1}{85}$ we also get the same thing so it's all good

 



\textbf{Solution 6 (Clear)}

Let $f(x)=\sum_{k=1}^{119}\lvert kx-1\rvert$. Note that the sum of functions that are concave up, as another function, is also concave up. Evaluating $f^\prime(x)$, we have
 \[f'(x)=\sum_{kx> 1}k -\sum_{kx<1}k\]Let $m=\left\lfloor \dfrac 1x\right\rfloor$. This means $m$ is the greatest integer that satisfies $m < \dfrac 1x \iff mx < 1$. So
 \[f'(x)= \sum_{k=m+1}^{119}k -\sum_{k=1}^m k\]We desire the integer value $m$ such that the derivative $f'(x)$ evaluates to $0$, or
 \[\frac{(119-m)(m+120)}{2}=\frac{m(m+1)}{2}\]Solving yields $m=84$. In the case of minimal of $f$, we have all $x$ satisfying $m\le \dfrac 1x < m+1\iff \frac{1}{85} < x \le \frac{1}{84}$ yields the minimal. For convenience, we evaluate the minimal of $f$ when $x=\frac{1}{84}$:\begin{align*}f(x)&=\sum_{k=1}^{84}\left(1-\frac{k}{84}\right)+\sum_{k=85}^{119}\left(\frac{k}{84}-1\right)\\&=84-\frac{85}{2}+\frac{85}{2}-35\\&=49\end{align*}
This leads to the option $\boxed{\textbf{(A)}\ 49}$.

 


	\item (\textit{AMC 12 12A 2010}) The number obtained from the last two nonzero digits of $90!$ is equal to $n$. What is $n$?


 $\textbf{(A)}\ 12 \qquad \textbf{(B)}\ 32 \qquad \textbf{(C)}\ 48 \qquad \textbf{(D)}\ 52 \qquad \textbf{(E)}\ 68$

 \textbf{Answer: }A



\textbf{Solution 1}

We will use the fact that for any integer $n$,
 \begin{align*}(5n+1)(5n+2)(5n+3)(5n+4)&=[(5n+4)(5n+1)][(5n+2)(5n+3)]\\ &=(25n^2+25n+4)(25n^2+25n+6)\equiv 4\cdot 6\\ &=24\pmod{25}\equiv -1\pmod{25}.\end{align*}


 First, we find that the number of factors of $10$ in $90!$ is equal to $\left\lfloor \frac{90}5\right\rfloor+\left\lfloor\frac{90}{25}\right\rfloor=18+3=21$. Let $N=\frac{90!}{10^{21}}$. The $n$ we want is therefore the last two digits of $N$, or $N\pmod{100}$. If instead we find $N\pmod{25}$, we know that $N\pmod{100}$, what we are looking for, could be $N\pmod{25}$, $N\pmod{25}+25$, $N\pmod{25}+50$, or $N\pmod{25}+75$. Only one of these numbers will be a multiple of four, and whichever one that is will be the answer, because $N\pmod{100}$ has to be a multiple of 4.

 If we divide $90!$ by $5^{21}$ to create $M$ by taking out all the factors of $5$ in $90!$, we can write $N$ as $\frac M{2^{21}}$ where
 \[M=1\cdot 2\cdot 3\cdot 4\cdot 1\cdot 6\cdot 7\cdot 8\cdot 9\cdot 2\cdots 89\cdot 18,\]where every multiple of 5 is replaced by the number with all its factors of 5 removed. Specifically, every number in the form $5n$ is replaced by $n$, and every number in the form $25n$ is replaced by $n$.

 The number $M$ can be grouped as follows:

 
 \begin{align*}M= &(1\cdot 2\cdot 3\cdot 4)(6\cdot 7\cdot 8\cdot 9)\cdots(86\cdot 87\cdot 88\cdot 89)\\ &\cdot (1\cdot 2\cdot 3\cdot 4)(6\cdot 7\cdot 8\cdot 9)\cdots (16\cdot 17\cdot 18) \\ &\cdot (1\cdot 2\cdot 3).\end{align*}


 Where the first line is composed of the numbers in $90!$ that aren't multiples of five, the second line is the multiples of five \textbf{and not 25} after they have been divided by five, and the third line is multiples of 25 after they have been divided by 25.

 Using the identity at the beginning of the solution, we can reduce $M$ to

 
 \begin{align*}M&\equiv(-1)^{18} \cdot (-1)^3(16\cdot 17\cdot 18) \cdot (1\cdot 2\cdot 3) \\ &= 1\cdot -21\cdot 6\\ &= -1\pmod{25} =24\pmod{25}.\end{align*}


 Using the fact that $2^{10}=1024\equiv -1\pmod{25}$ (or simply the fact that $2^{21}=2097152$ if you have your powers of 2 memorized), we can deduce that $2^{21}\equiv 2\pmod{25}$. Therefore $N=\frac M{2^{21}}\equiv \frac {24}2\pmod{25}=12\pmod{25}$.

 Finally, combining with the fact that $N\equiv 0\pmod 4$ yields $n=\boxed{\textbf{(A)}\ 12}$.

 



\textbf{Solution 2}

Let $P$ be $90!$ after we truncate its zeros. Notice that $90!$ has exactly (floored) $\left\lfloor\frac{90}{5}\right\rfloor + \left\lfloor\frac{90}{25}\right\rfloor = 21$ factors of 5; thus, 
 \[P = 2^{-21}\cdot 5^{-21}\cdot 90!.\] We shall consider $P$ modulo 4 and 25, to determine its residue modulo 100. It is easy to prove that $P$ is divisible by 4 (consider the number of 2s dividing $90!$ minus the number of 5s dividing $90!$), and so we only need to consider $P$ modulo 25.

 Now, notice that for integers $a, n$ we have
 \[(5n + a)(5n - a) \equiv -a^2 \mod 25.\]

 Thus, for integral a: 
 \[(10a + 1)(10a + 2)(10a + 3)(10a + 4)(10a + 6)(10a + 7)(10a + 8)(10a + 9) \equiv (-1)(-4)(-9)(-16) \equiv 576 \equiv 1 \mod 25.\] Using this process, we can essentially remove all the numbers which had not formerly been a multiple of 5 in $90!$ from consideration.

 Now, we consider the remnants of the 5, 10, 15, 20, ..., 90 not yet eliminated. After factoring out a 2 from each of the 9 even numbers in this sequence, the 10, 20, 30, ..., 90 becomes 1, 2, 3, 4, 1, 6, 7, 8, 9, whose product is 1 mod 25. We have accounted for 9 of the 21 2's thus we still need to multiply by $2^{-12}$. Also, the 5, 15, 25, ..., 85 becomes 1, 3, 1, 7, 9, 11, 13, 3, 17. Multiplying out 1, 3, 1, 7, ..., 17 yields 2 modulo 25, and so we are left to compute $2^{-11}$. Also note that $\phi(25)=5^2-5^1=20$. Euler's Theorem states that $2^{\phi(25)} \equiv 1 \mod 25$, thus $2^{-11} \equiv 2^{-11+20} = 2^9 = 512 \equiv 12 \mod 25$. Because 12 is also a multiple of 4, we can utilize the Chinese Remainder Theorem to show that $P \equiv 12 \mod 100$ and so the answer is $\boxed{12}$.

 



\textbf{Solution 3}

We start of by truncating the $0$s off $90!$, just like Solution 2. Since there are $v_5(90!) = \left\lfloor\frac{90}{5}\right\rfloor+\left\lfloor\frac{90}{25}\right\rfloor = 21$ terminating zeroes, we have the number we obtain from truncating the terminating zeroes at the end of $90!$ will be $\frac{90!}{10^{21}}$.

 By \href{https://artofproblemsolving.com/wiki/index.php?title=Chinese\_Remainder\_Theorem}{Chinese Remainder Theorem}, we can divide the mod $100$ into mod $4$ and mod $25$. We know that there are way more than $23$ $2$s in $90!$, so we have

 
 \[\dfrac{90!}{10^{21}} \equiv 0 \text{ } \text{ (mod } 4 \text{)}\]

 Now, notice that $\frac{90!}{5^{21}}$ is basically $90!$ without any factors that are multiples of $5$.

 
 \[1\cdot 2\cdot 3\cdot 4\cdot 6\cdot 7\cdots 23\cdot 24 \text{ } \text{ (mod } 25 \text{)}\]

 We can rewrite the expression as 

 
 \[(1\cdot 2\cdot 3\cdot 4\cdot 6\cdot 7\cdots 12)^2 \text{ } \text{ (mod } 25 \text{)}\]

 We also have 

 
 \[1\cdot 2\cdot 3\cdot 4\cdot 6\cdot 7\cdots 12 \equiv 8\cdot 8\cdot 9\cdot 11\cdot 12 \equiv 7 \text{ } \text{ (mod } 25 \text{)}\]

 Squaring the expression gets us

 
 \[(1\cdot 2\cdot 3\cdot 4\cdot 6\cdot 7\cdots 12)^2 \equiv 49 \equiv -1 \text{ } \text{ (mod } 25 \text{)}\]

 Notice that there are $3$ of this string of numbers multiplied together occurs four times in $\frac{90!}{10^{21}}$, the fourth being only a partial

 $1\cdot 2\cdot 3\cdot 4\cdots 14 \equiv 24 \equiv -1 \text{ } \text{ (mod } 25 \text{)}$

 Note that as in Solution 1, the residue from extracting the fives is$1\cdot 2\cdot 3\cdot 4\cdots 18 \cdot 1 \cdot 2 \cdot 3 \equiv -1 \text{ } \text{ (mod } 25 \text{)}$

 Multiplying these together gives us

 
 \[\frac{90!}{10^{21}} \equiv \frac{\frac{90!}{5^{21}}}{2^{21}} \equiv \frac{(-1)^5}{2^{21}} \equiv \frac{-1}{2^{21}} \text{ } \text{ (mod } 25 \text{)}\]

 By Euler's Theorem, we have $\frac{-1}{2^{21}} \equiv -13^{21} \equiv -13^{\phi{(25)}}\cdot 13 \equiv -13 \equiv 12 \text{ } \text{ (mod } 25 \text{)}$

 Thus, we want to solve $x \equiv 12 \text{ } \text{ (mod } 25 \text{)}$ and $x \equiv 0 \text{ } \text{ (mod } 4 \text{)}$. An obvious solution is $x\equiv 12$ and by Chinese Remainder Theorem, we have $x\equiv 12$ is the only solution. So thus, $\frac{90!}{10^{21}} \equiv \boxed{12} \text{ } \text{ (mod } 100 \text{)}$

 \textbf{Remark (Chinese Remainder Theorem)} Both solution 1 and solution 2 rely on $n \equiv 12 (\bmod { \quad 25})$, $n \equiv 0 (\bmod { \quad 4})$ to get $n (\bmod { \quad 100})$

 By \href{https://artofproblemsolving.com/wiki/index.php/Chinese\_Remainder\_Theorem}{Chinese Remainder Theorem}, the general solution of the system of $2$ linear congruences is:

 \begin{verbatim}
$n \equiv r_1 (\bmod { \quad m_1})$, $n \equiv r_2 (\bmod { \quad m_2})$, $(m_1, m_2) = 1$Find $k_1$ and $k_2$ such that $k_1 m_1 \equiv 1 (\bmod{ \quad m_2})$, $k_2 m_2 \equiv 1 (\bmod{ \quad m_1})$Then $n \equiv k_2 m_2 r_1 + k_1 m_1 r_2 (\bmod{ \quad m_1 m_2})$
\end{verbatim}

In this problem, $n \equiv 12 (\bmod { \quad 25})$, $n \equiv 0 (\bmod { \quad 4})$:

 \begin{verbatim}
$n \equiv 12 (\bmod{ \quad 25})$, $n \equiv 0 (\bmod{ \quad 4})$, $(25, 4) = 1$$1 \cdot 25 \equiv 1 (\bmod{ \quad 4})$, $19 \cdot 4 \equiv 1 (\bmod{ \quad 25})$Then $n \equiv 19 \cdot 4 \cdot 12 + 1 \cdot 25 \cdot 0 \equiv 12 (\bmod{ \quad 100})$
\end{verbatim}




	\item (\textit{AMC 12 12A 2010}) Let $f(x) = \log_{10} \left(\sin(\pi x) \cdot \sin(2 \pi x) \cdot \sin (3 \pi x) \cdots \sin(8 \pi x)\right)$. The intersection of the domain of $f(x)$ with the interval $[0,1]$ is a union of $n$ disjoint open intervals. What is $n$?


 $\textbf{(A)}\ 2 \qquad \textbf{(B)}\ 12 \qquad \textbf{(C)}\ 18 \qquad \textbf{(D)}\ 22 \qquad \textbf{(E)}\ 36$

 \textbf{Answer: }B



\textbf{Solution 1}

The question asks for the number of disjoint open intervals, which means we need to find the number of disjoint intervals such that the function is defined within them.

 We note that since all of the $\sin$ factors are inside a logarithm, the function is undefined where the inside of the logarithm is less than or equal to $0$.

 First, let us find the number of zeros of the inside of the logarithm.

 
 \begin{align*}\sin(\pi x) \cdot \sin(2 \pi x) \cdot \sin (3 \pi x) \cdots \sin(8 \pi x) &= 0\\ \sin(\pi x) &= 0\\ x &= 0, 1\\ \sin(2 \pi x) &= 0\\ x &= 0, \frac{1}{2}, 1\\ \sin(3 \pi x) &= 0\\ x &= 0, \frac{1}{3}, \frac{2}{3}, 1\\ \sin(4 \pi x) &= 0\\ x &= 0, \frac{1}{4}, \frac{2}{4}, \frac{3}{4}, 1\\ &\cdots\end{align*}


 After counting up the number of zeros for each factor and eliminating the excess cases we get $23$ zeros and $22$ intervals.

 In order to find which intervals are negative, we must first realize that at every zero of each factor, the sign changes. We also have to be careful, as some zeros are doubled, or even tripled, quadrupled, etc.

 The first interval $\left(0, \frac{1}{8}\right)$ is obviously positive. This means the next interval $\left(\frac{1}{8}, \frac{1}{7}\right)$ is negative. Continuing the pattern and accounting for doubled roots (which do not flip sign), we realize that there are 

 


	\item (\textit{AMC 12 12A 2010}) Two quadrilaterals are considered the same if one can be obtained from the other by a rotation and a translation. How many different convex cyclic quadrilaterals are there with integer sides and perimeter equal to 32?


 $\textbf{(A)}\ 560 \qquad \textbf{(B)}\ 564 \qquad \textbf{(C)}\ 568 \qquad \textbf{(D)}\ 1498 \qquad \textbf{(E)}\ 2255$

 \textbf{Answer: }C



\textbf{Solution 1}

It should first be noted that given any quadrilateral of fixed side lengths, there is exactly one way to manipulate the angles so that the quadrilateral becomes cyclic.

 \textbf{Proof.} Given a quadrilateral $ABCD$ where all sides are fixed (in a certain order), we can construct the diagonal $\overline{BD}$. When $BD$ is the minimum allowed by the triangle inequality, one of the angles $\angle DAB$ or $\angle BCD$ will be degenerate and measure $0^\circ$, so opposite angles will sum to less than $180^\circ$. When $BD$ is the maximum allowed, one of the angles will be degenerate and measure $180^\circ$, so opposite angles will sum to more than $180^\circ$. Thus, since the sum of opposite angles increases continuously as $BD$ is lengthened from the minimum to the maximum values, there is a unique value of $BD$ somewhere in the middle such that the sum of opposite angles is exactly $180^\circ$ (due to IVT).

 Denote $a$, $b$, $c$, and $d$ as the integer side lengths of the quadrilateral. Without loss of generality, let $a\ge b \ge c \ge d$.

 Since $a+b+c+d = 32$, the \href{https://artofproblemsolving.com/wiki/index.php?title=Triangle\_Inequality}{Triangle Inequality} implies that $a \le 15$.

 
We will now split into $5$ cases.

 
Case $1$: $a = b = c = d$ ($4$ side lengths are equal)

 Clearly there is only $1$ way to select the side lengths $(8,8,8,8)$, and no matter how the sides are rearranged only $1$ unique quadrilateral can be formed.

 Case $2$: $a = b = c > d$ or $a > b = c = d$ ($3$ side lengths are equal)

 If $3$ side lengths are equal, then each of those side lengths can only be integers from $6$ to $10$ except for $8$ (because that is counted in the first case). Obviously there is still only $1$ unique quadrilateral that can be formed from one set of side lengths, resulting in a total of $4$ quadrilaterals.

 Case $3$: $a = b > c = d$ ($2$ pairs of side lengths are equal)

 $a$ and $b$ can be any integer from $9$ to $15$, and likewise $c$ and $d$ can be any integer from $1$ to $7$. However, a single set of side lengths can form $2$ different cyclic quadrilaterals (a rectangle and a kite), so the total number of quadrilaterals for this case is $7\cdot{2} = 14$.

 Case $4$: $a = b > c > d$ or $a > b = c > d$ or $a > b > c = d$ ($2$ side lengths are equal)

 If the $2$ equal side lengths are each $1$, then the other $2$ sides must each be $15$, which we have already counted in an earlier case. If the equal side lengths are each $2$, there is $1$ possible set of side lengths. Likewise, for side lengths of $3$ there are $2$ sets. Continuing this pattern, we find a total of $1+2+3+4+4+5+7+5+4+4+3+2+1 = 45$ sets of side lengths. (Be VERY careful when adding up the total for this case!) For each set of side lengths, there are $3$ possible quadrilaterals that can be formed, so the total number of quadrilaterals for this case is $3\cdot{45} = 135$.

 Case $5$: $a > b > c > d$ (no side lengths are equal)Using the same counting principles starting from $a = 15$ and eventually reaching $a = 9$, we find that the total number of possible side lengths is $69$. There are $4!$ ways to arrange the $4$ side lengths, but there is only $1$ unique quadrilateral for $4$ rotations, so the number of quadrilaterals for each set of side lengths is $\frac{4!}{4} = 6$. The total number of quadrilaterals is $6\cdot{69} = 414$.

 
And so, the total number of quadrilaterals that can be made is $414 + 135 + 14 + 4 + 1 = \boxed{568\ \textbf{(C)}}$.

 



\textbf{Solution 2}

As with solution $1$ we would  like to note that given any quadrilateral we can change its angles to make a cyclic one.

 Let $a \ge b \ge c\ge d$ be the sides of the quadrilateral.

 There are $\binom{31}{3}$ ways to partition $32$. However, some of these will not be quadrilaterals since they would have one side bigger than the sum of the other three.  This occurs when $a \ge 16$.  For $a=16$, $b+c+d=16$. There are $\binom{15}{2}$ ways to partition $16$. Since $a$ could be any of the four sides, we have counted $4\binom{15}{2}$ degenerate quadrilaterals. Similarly, there are $4\binom{14}{2}$, $4\binom{13}{2} \cdots 4\binom{2}{2}$ for other values of $a$.  Thus, there are $\binom{31}{3} - 4\left(\binom{15}{2}+\binom{14}{2}+\cdots+\binom{2}{2}\right) = \binom{31}{3} - 4\binom{16}{3} = 2255$ non-degenerate partitions of $32$ by the hockey stick theorem.  We then account for symmetry. If all sides are congruent (meaning the quadrilateral is a square), the quadrilateral will be counted once. If the quadrilateral is a rectangle (and not a square), it will be counted twice. In all other cases, it will be counted 4 times. Since there is $1$ square case, and $7$ rectangle cases, there are $2255-1-2\cdot7=2240$ quadrilaterals counted 4 times. Thus there are $1+7+\frac{2240}{4} = \boxed{568} = \boxed{\textbf{(C)}}$ total quadrilaterals.

 



\textbf{Solution 3 (Burnside)}

As with solution $2$ we find that there are $2255$ ways to form a quadrilateral if we don't account for rotations. We now apply \href{https://artofproblemsolving.com/wiki/index.php?title=Burnside\%27s\_Lemma}{Burnside's Lemma}. There are four types of actions in the group acting on the set of quadrilaterals. We will consider each individually:Identity: maps a quadrilateral with sides $a,b,c,d$ in that order to $a,b,c,d$. Obviously all members of the set of quadrilaterals are fixed points, for a total of $2255$.Rotation by one: maps a quadrilateral from $a,b,c,d$ to $b,c,d,a$. For this to have a fixed point we need $a=b,b=c,c=d,d=a$, so the only quadrilateral that is a fixed point is the square with side length $8$, for a total of $1$.Rotation by two: maps a quadrilateral from $a,b,c,d$ to $c,d,a,b$. For this to be a fixed point we need $a=c$ and $b=d$. Thus the quadrilateral is of the form $x,y,x,y$—a rectangle. We can count that there are $15$ rectangle cases, namely $(a, b, c, d) = \{ (1, 15, 1, 15), (2, 14, 2, 14), \cdots, (15, 1, 15, 1)  \}$.Rotation by three: maps a quadrilateral from $a,b,c,d$ to $d,a,b,c$. Similarly to the rotation by one case, there is one fixed point here.Summing up, we get that the total number of groups is $2255+1+15+1=2272$. Since there are $4$ members of the group our final answer is $\frac{2272}{4}=\boxed{568}$ total orbits of the set of quadrilaterals, so the answer is $\boxed{\textbf{C}}$.

 


\end{enumerate}